
\textbf{Yiqiao Liu, final-week Ed and oral feedback (author-confirmed October 21).} Yiqiao Liu introduced Armory and challenged novelty beyond scheduling plus critical payments. \emph{Revise}: Sections 1 and 7, Figure 1 and Appendix C now compare RTC, Armory and Themis and limit novelty. \emph{Keep}: the physical allocation/payments connection survives. \emph{Reasoned disagreement}: throughput alone cannot assess private-value incentives; add realized value and entrant access, reporting payments separately. \emph{Investigate}: a matched dynamic-batching benchmark requires a common execution model; no Armory superiority is claimed. The author supplied the comment and confirmed its attribution; Ed record: \url{https://edstem.org/us/courses/102917/discussion/8381683?comment=19460483}. Responsible author: Xinrong Sun.

\textbf{October 14 author self-challenge.} Could a false deadline improve utility? \emph{Test and revise}, October 20: 88,200 finite deviations found no gain; Appendix G supplies the conditional argument and a service-duration counterexample with gain 8. This changes the claim from general suspicion of deadline reports to a specific verification boundary. Locations: Section 6, Appendix G, code/deadline\_test.py and the demonstration's self-test panel. Responsible author: Xinrong Sun.

\textbf{Instructor final rubric, supplied October 20.} \emph{Revise}: migrate five PS2 sections into the eight required parts; retain A--F and add G; create the impact and roadmap diagrams, official SDG icon, contribution and credit blocks, and repository poster preview. Verification: source compilation, page-count and poster-dimension checks, fresh-run outputs and package audit. Responsible author: Xinrong Sun.


\textbf{Xuantong Fu (FP3), supplied peer-feedback record.} The false-label example gives $A$ a gain of 4, challenging unconditional truthfulness; code results should reproduce across artifacts. \emph{Revise}: distinguish value reports from eligibility reports and state verified eligibility as a condition. The legacy profile $(10,9,6)$, reserve 1 and credit 5 changes $A$'s utility from 0 to 4 when it falsely claims eligibility. This is the earlier unit-window model, not the final duration test. \emph{Keep}: conditional value truthfulness; \emph{Investigate}: qualification auditing. Verification: rerun the legacy counterexample, retain robot and JavaScript parity checks, and identify the final release. Locations: Appendix A, \path{code/model.py}, \path{evidence/final-release.json}. Responsible author: Xinrong Sun.

\textbf{Han Zhang, discussion reported by the author.} Distinguish robot inference from an ordinary CPU auction. \emph{Revise}: action-buffer expiry and target-pose validity determine physical deadlines; relaxing either constraint changes the worked allocation. Verification: the two state ablations in \path{code/robot_results.json}. Locations: Sections 1, 3, 5--6 and Appendix A. Responsible author: Xinrong Sun.

\textbf{Prof. Ken Rogerson and Muhan Chen.} The author reports that Prof. Rogerson offered mechanism-improvement feedback, while Muhan Chen supplied oral and handwritten suggestions. Their contributions are credited without attributing an unrecorded specific proposition to them. The final mechanism, limits and response record make the resulting project available for their review. The author's acknowledgment post is \url{https://edstem.org/us/courses/102917/discussion/8381683?comment=19460597}.
